\clearpage
\appendix
\section{小さなパラメータでの分類表}
\label{app:tables}

$n$ は変数数、$\dim R=n-1$、$a=h-\sum_iw_i$ とする。$3\leq n\leq5$、$1\leq a\leq6$ の全18入力を Lean 列挙器で実行した。表の数値は次数付き複素代数同型類の数である。

\subsection{個数と渡辺氏のリストとの比較}
\begin{table}[htbp]
\centering
% Generated by scripts/blueprint_tables.py; do not edit by hand.
\begin{tabular}{ccrrrrrr}
\hline
$n$ & $\dim R$ & $a=1$ & $a=2$ & $a=3$ & $a=4$ & $a=5$ & $a=6$ \\
\hline
3 & 2 & 14 & 6 & 8 & 8 & 21 & 0 \\
4 & 3 & 1 & 0 & 0 & 0 & 1 & 0 \\
5 & 4 & 3 & 0 & 0 & 0 & 2 & 0 \\
\hline
\end{tabular}

\caption{小さな $n,a$ に対する同型類の全個数。}
\label{tab:counts}
\end{table}

表\ref{tab:comparison}の $N_{\mathrm W}$ は渡辺氏の arXiv 第1版\cite{Watanabe}の種数零・三点型の項目数、$N_{\mathrm{Lean}}$ は $n=3$ の出力数である。比較項目は (1-C-1)--(1-C-14)、(2-C-4)--(2-C-9)、(3-C-14)--(3-C-20)、(4-C-6)--(4-C-12)、(5-C-26)--(5-C-31)、(5-C-34)--(5-C-44)。$a=6$ の該当項目はない。(3-C-19) の $(4,10.17;34)$ は $(4,10,17;34)$ と読み、欠落に数えない。

\begin{table}[htbp]
\centering
% Generated by scripts/blueprint_tables.py; do not edit by hand.
\begin{tabular}{crrr}
\hline
$a$ & $N_{\mathrm{W}}$ & $N_{\mathrm{Lean}}$ & $N_{\mathrm{missing}}$ \\
\hline
1 & 14 & 14 & 0 \\
2 & 6 & 6 & 0 \\
3 & 7 & 8 & 1 \\
4 & 7 & 8 & 1 \\
5 & 17 & 21 & 4 \\
6 & 0 & 0 & 0 \\
\hline
\end{tabular}

\caption{\texttt{arXiv:1401.0789v1} の指定した三点型項目との比較。}
\label{tab:comparison}
\end{table}

転記した51項目は全て列挙器の57同型類に含まれる。表\ref{tab:missing}が追加の六例である。$(x,y,z)$ の重みは表示順、signature は独立に整列する。各式は三項で、重みは原始的、$|\det E|=h$ である。

\begin{table}[htbp]
\centering\small
% Generated by scripts/blueprint_tables.py; do not edit by hand.
\begin{tabular}{cclc}
\hline
$a$ & $(w_1,w_2,w_3;h)$ & $f(x,y,z)$ & $(p_1,p_2,p_3)$ \\
\hline
3 & $(4,10,13;30)$ & $xz^{2}+x^{5}y+y^{3}$ & $(2,4,13)$ \\
4 & $(5,6,9;24)$ & $x^{3}z+yz^{2}+y^{4}$ & $(3,5,9)$ \\
5 & $(4,7,9;25)$ & $yz^{2}+xy^{3}+x^{4}z$ & $(4,7,9)$ \\
5 & $(4,7,12;28)$ & $xz^{2}+x^{4}z+y^{4}$ & $(4,4,12)$ \\
5 & $(4,7,16;32)$ & $xy^{4}+x^{4}z+z^{2}$ & $(4,4,7)$ \\
5 & $(6,8,19;38)$ & $xy^{4}+x^{5}y+z^{2}$ & $(2,6,8)$ \\
\hline
\end{tabular}

\caption{比較対象のリストに載っていない六つの重み系。}
\label{tab:missing}
\end{table}

\clearpage
\subsection{小さなパラメータでの三変数の完全なリスト}
表\ref{tab:ternary1}--\ref{tab:ternary5}に $a=1,\ldots,5$ の全同型類を示す。$a=6$ は空である。星印は表\ref{tab:missing}の補完例を表し、それ以外のキーは比較対象のリストに含まれる。$a=1$ の表は十四個の例外型単峰超曲面三角特異点を再現する。

\begin{table}[htbp]
\centering\small
% Generated by scripts/blueprint_tables.py; do not edit by hand.
\begin{tabular}{clcc}
\hline
$(w_1,w_2,w_3;h)$ & $f(x,y,z)$ & $(p_1,p_2,p_3)$ &  \\
\hline
$(3,4,4;12)$ & $y^{2}z+yz^{2}+x^{4}$ & $(4,4,4)$ &  \\
$(3,4,5;13)$ & $y^{2}z+xz^{2}+x^{3}y$ & $(3,4,5)$ &  \\
$(3,4,8;16)$ & $x^{4}y+y^{2}z+z^{2}$ & $(3,4,4)$ &  \\
$(3,5,6;15)$ & $xz^{2}+x^{3}z+y^{3}$ & $(3,3,6)$ &  \\
$(3,5,9;18)$ & $xy^{3}+x^{3}z+z^{2}$ & $(3,3,5)$ &  \\
$(3,8,12;24)$ & $x^{4}z+z^{2}+y^{3}$ & $(3,3,4)$ &  \\
$(4,5,6;16)$ & $y^{2}z+xz^{2}+x^{4}$ & $(2,5,6)$ &  \\
$(4,5,10;20)$ & $y^{2}z+z^{2}+x^{5}$ & $(2,5,5)$ &  \\
$(4,6,7;18)$ & $xz^{2}+x^{3}y+y^{3}$ & $(2,4,7)$ &  \\
$(4,6,11;22)$ & $xy^{3}+x^{4}y+z^{2}$ & $(2,4,6)$ &  \\
$(4,10,15;30)$ & $x^{5}y+y^{3}+z^{2}$ & $(2,4,5)$ &  \\
$(6,8,9;24)$ & $xz^{2}+x^{4}+y^{3}$ & $(2,3,9)$ &  \\
$(6,8,15;30)$ & $xy^{3}+x^{5}+z^{2}$ & $(2,3,8)$ &  \\
$(6,14,21;42)$ & $z^{2}+y^{3}+x^{7}$ & $(2,3,7)$ &  \\
\hline
\end{tabular}

\caption{$n=3$, $a=1$：14同型類。}
\label{tab:ternary1}
\end{table}
\begin{table}[htbp]
\centering\small
% Generated by scripts/blueprint_tables.py; do not edit by hand.
\begin{tabular}{clcc}
\hline
$(w_1,w_2,w_3;h)$ & $f(x,y,z)$ & $(p_1,p_2,p_3)$ &  \\
\hline
$(3,5,5;15)$ & $y^{2}z+yz^{2}+x^{5}$ & $(5,5,5)$ &  \\
$(3,5,7;17)$ & $y^{2}z+xz^{2}+x^{4}y$ & $(3,5,7)$ &  \\
$(3,5,10;20)$ & $x^{5}y+y^{2}z+z^{2}$ & $(3,5,5)$ &  \\
$(3,7,9;21)$ & $xz^{2}+x^{4}z+y^{3}$ & $(3,3,9)$ &  \\
$(3,7,12;24)$ & $xy^{3}+x^{4}z+z^{2}$ & $(3,3,7)$ &  \\
$(3,10,15;30)$ & $x^{5}z+z^{2}+y^{3}$ & $(3,3,5)$ &  \\
\hline
\end{tabular}

\caption{$n=3$, $a=2$：6同型類。}
\label{tab:ternary2}
\end{table}
\begin{table}[htbp]
\centering\small
% Generated by scripts/blueprint_tables.py; do not edit by hand.
\begin{tabular}{clcc}
\hline
$(w_1,w_2,w_3;h)$ & $f(x,y,z)$ & $(p_1,p_2,p_3)$ &  \\
\hline
$(4,5,7;19)$ & $yz^{2}+xy^{3}+x^{3}z$ & $(4,5,7)$ &  \\
$(4,5,8;20)$ & $xz^{2}+x^{3}z+y^{4}$ & $(4,4,8)$ &  \\
$(4,5,12;24)$ & $xy^{4}+x^{3}z+z^{2}$ & $(4,4,5)$ &  \\
$(4,7,10;24)$ & $y^{2}z+xz^{2}+x^{6}$ & $(2,7,10)$ &  \\
$(4,7,14;28)$ & $y^{2}z+z^{2}+x^{7}$ & $(2,7,7)$ &  \\
$(4,10,13;30)$ & $xz^{2}+x^{5}y+y^{3}$ & $(2,4,13)$ & $\ast$ \\
$(4,10,17;34)$ & $xy^{3}+x^{6}y+z^{2}$ & $(2,4,10)$ &  \\
$(4,14,21;42)$ & $x^{7}y+y^{3}+z^{2}$ & $(2,4,7)$ &  \\
\hline
\end{tabular}

\caption{$n=3$, $a=3$：8同型類。}
\label{tab:ternary3}
\end{table}
\begin{table}[htbp]
\centering\small
% Generated by scripts/blueprint_tables.py; do not edit by hand.
\begin{tabular}{clcc}
\hline
$(w_1,w_2,w_3;h)$ & $f(x,y,z)$ & $(p_1,p_2,p_3)$ &  \\
\hline
$(3,7,7;21)$ & $y^{2}z+yz^{2}+x^{7}$ & $(7,7,7)$ &  \\
$(3,7,11;25)$ & $y^{2}z+xz^{2}+x^{6}y$ & $(3,7,11)$ &  \\
$(3,7,14;28)$ & $x^{7}y+y^{2}z+z^{2}$ & $(3,7,7)$ &  \\
$(3,11,15;33)$ & $xz^{2}+x^{6}z+y^{3}$ & $(3,3,15)$ &  \\
$(3,11,18;36)$ & $xy^{3}+x^{6}z+z^{2}$ & $(3,3,11)$ &  \\
$(3,14,21;42)$ & $x^{7}z+z^{2}+y^{3}$ & $(3,3,7)$ &  \\
$(5,6,9;24)$ & $x^{3}z+yz^{2}+y^{4}$ & $(3,5,9)$ & $\ast$ \\
$(5,6,15;30)$ & $x^{3}z+z^{2}+y^{5}$ & $(3,5,5)$ &  \\
\hline
\end{tabular}

\caption{$n=3$, $a=4$：8同型類。}
\label{tab:ternary4}
\end{table}
\begin{table}[htbp]
\centering\small
% Generated by scripts/blueprint_tables.py; do not edit by hand.
\begin{tabular}{clcc}
\hline
$(w_1,w_2,w_3;h)$ & $f(x,y,z)$ & $(p_1,p_2,p_3)$ &  \\
\hline
$(3,8,8;24)$ & $y^{2}z+yz^{2}+x^{8}$ & $(8,8,8)$ &  \\
$(3,8,13;29)$ & $y^{2}z+xz^{2}+x^{7}y$ & $(3,8,13)$ &  \\
$(3,8,16;32)$ & $x^{8}y+y^{2}z+z^{2}$ & $(3,8,8)$ &  \\
$(3,13,18;39)$ & $xz^{2}+x^{7}z+y^{3}$ & $(3,3,18)$ &  \\
$(3,13,21;42)$ & $xy^{3}+x^{7}z+z^{2}$ & $(3,3,13)$ &  \\
$(3,16,24;48)$ & $x^{8}z+z^{2}+y^{3}$ & $(3,3,8)$ &  \\
$(4,7,9;25)$ & $yz^{2}+xy^{3}+x^{4}z$ & $(4,7,9)$ & $\ast$ \\
$(4,7,12;28)$ & $xz^{2}+x^{4}z+y^{4}$ & $(4,4,12)$ & $\ast$ \\
$(4,7,16;32)$ & $xy^{4}+x^{4}z+z^{2}$ & $(4,4,7)$ & $\ast$ \\
$(4,9,14;32)$ & $y^{2}z+xz^{2}+x^{8}$ & $(2,9,14)$ &  \\
$(4,9,18;36)$ & $y^{2}z+z^{2}+x^{9}$ & $(2,9,9)$ &  \\
$(4,14,19;42)$ & $xz^{2}+x^{7}y+y^{3}$ & $(2,4,19)$ &  \\
$(4,14,23;46)$ & $xy^{3}+x^{8}y+z^{2}$ & $(2,4,14)$ &  \\
$(4,18,27;54)$ & $x^{9}y+y^{3}+z^{2}$ & $(2,4,9)$ &  \\
$(6,7,9;27)$ & $xy^{3}+x^{3}z+z^{3}$ & $(3,6,7)$ &  \\
$(6,8,11;30)$ & $yz^{2}+xy^{3}+x^{5}$ & $(2,8,11)$ &  \\
$(6,8,13;32)$ & $xz^{2}+x^{4}y+y^{4}$ & $(2,6,13)$ &  \\
$(6,8,19;38)$ & $xy^{4}+x^{5}y+z^{2}$ & $(2,6,8)$ & $\ast$ \\
$(6,16,21;48)$ & $xz^{2}+x^{8}+y^{3}$ & $(2,3,21)$ &  \\
$(6,16,27;54)$ & $xy^{3}+x^{9}+z^{2}$ & $(2,3,16)$ &  \\
$(6,22,33;66)$ & $z^{2}+y^{3}+x^{11}$ & $(2,3,11)$ &  \\
\hline
\end{tabular}

\caption{$n=3$, $a=5$：21同型類。}
\label{tab:ternary5}
\end{table}

\clearpage
\subsection{小さなパラメータでの高次元の完全なリスト}
$n\geq4$ では各行の Fermat 表示は $f=\sum_i z_i^{h/w_i}$ である。signature と重みはそれぞれ独立に昇順で表示するので、変数の指数は重みの順序に従う。この範囲で $n=4,5$ の非空な結果があるのは $a=1,5$ だけである。

\begin{table}[htbp]
\centering\small
% Generated by scripts/blueprint_tables.py; do not edit by hand.
\begin{tabular}{ccc}
\hline
$a$ & $(p_1,\ldots,p_n)$ & $(w_1,\ldots,w_n;h)$ \\
\hline
1 & $(2,3,7,43)$ & $(42,258,602,903;1806)$ \\
5 & $(2,3,7,47)$ & $(42,282,658,987;1974)$ \\
\hline
\end{tabular}

\caption{$n=4$（次元3）、$1\leq a\leq6$ の全同型類。}
\label{tab:higher4}
\end{table}
\begin{table}[htbp]
\centering\small
% Generated by scripts/blueprint_tables.py; do not edit by hand.
\begin{tabular}{ccc}
\hline
$a$ & $(p_1,\ldots,p_n)$ & $(w_1,\ldots,w_n;h)$ \\
\hline
1 & $(2,3,11,23,31)$ & $(1518,2046,4278,15686,23529;47058)$ \\
1 & $(2,3,7,43,1807)$ & $(1806,75894,466206,1087814,1631721;3263442)$ \\
1 & $(2,3,7,47,395)$ & $(1974,16590,111390,259910,389865;779730)$ \\
5 & $(2,3,7,43,1811)$ & $(1806,76062,467238,1090222,1635333;3270666)$ \\
5 & $(2,3,7,71,103)$ & $(2982,4326,43878,102382,153573;307146)$ \\
\hline
\end{tabular}

\caption{$n=5$（次元4）、$1\leq a\leq6$ の全同型類。}
\label{tab:higher5}
\end{table}

\subsection{表の再現方法}
英語正本の \texttt{formal/} で \texttt{lake build canonical\_roots} を実行した後、そのリポジトリ直下で \texttt{python3 scripts/blueprint\_tables.py --check} を実行する。18入力を再実行し、保存した JSON と TeX の表と比較する。日本語版の表は同じコマンドに \texttt{--output-root} と日本語版リポジトリへのパスを追加して確認できる。意図的に再生成するときは \texttt{--check} を外す。

数値データは \texttt{blueprint/data/small\_classifications.json} にある。原資料の項目番号とページを付けた転記は、英語正本の \texttt{research/current/watanabe\_v1\_three\_point\_keys.csv} に保存している。分類の普遍的な正しさの定理に加え、\texttt{OutputCertificates.lean} は $n=3$, $1\leq a\leq6$ と $(n,a)=(4,1),(4,5),(5,1)$ などの全リスト等式を Lean のカーネルで検証している。表は実行出力から生成して照合する。歴史的な文献の読解・転記の確認は、この形式証明の外側にある。
